
\documentclass[11pt,a4paper]{article}

\usepackage[margin=1in]{geometry}
\usepackage{amsmath,amssymb,amsthm,mathtools}
\usepackage{booktabs,array}
\usepackage{enumitem}
\usepackage{hyperref}
\usepackage{xcolor}
\usepackage{microtype}

\hypersetup{
  colorlinks=true,
  linkcolor=blue!55!black,
  citecolor=blue!55!black,
  urlcolor=blue!55!black
}

\setlist[itemize]{leftmargin=1.5em}
\setlist[enumerate]{leftmargin=1.8em}
\allowdisplaybreaks

\newtheorem{proposition}{Proposition}
\newtheorem{lemma}{Lemma}
\newtheorem{remark}{Remark}
\newtheorem{assumption}{Assumption}

\newcommand{\R}{\mathbb{R}}
\newcommand{\E}{\mathbb{E}}
\newcommand{\1}{\mathbf{1}}
\newcommand{\DeltaN}{\Delta_n}
\newcommand{\pos}[1]{(#1)_+}
\newcommand{\MSE}{\operatorname{MSE}}
\newcommand{\Lap}{\operatorname{Laplace}}
\newcommand{\argmin}{\operatorname*{arg\,min}}

\title{\textbf{Fairness-Augmented Differentially Private Data Acquisition}\\
}
\author{Working notes for research development}
\date{15 November 2026}

\begin{document}
\maketitle
\tableofcontents
\newpage

\section{Purpose and status of this document}

These notes consolidate the model and derivations developed while studying
\emph{Optimal and Differentially Private Data Acquisition: Central and Local
Mechanisms} (the uploaded source paper), and while extending its central
mechanism-design problem with a two-group fairness penalty.

The main technical objective is to understand whether the paper's efficient
threshold-sweep algorithm can be extended to the fairness-augmented problem.
The current derivation suggests that it can: although fairness creates two
group-specific thresholds, the consistency equation couples them along a
unique monotone, piecewise-linear path. The outer stationarity condition is
quartic on each active-set segment.

\paragraph{Status warning.}
This is developed by AI based on notes and resources supplied. No data or integrity violations intended. Please report any mistakes.

\section{Original data-acquisition model}

\subsection{Estimator and privacy allocation}

The central mechanism considered in the source paper uses a weighted linear
estimator with Laplace noise:
\begin{equation}
  \widehat{\theta}
  = \sum_{i=1}^n w_i X_i + Z,
  \qquad
  Z\sim \Lap(1/\eta),
  \label{eq:estimator}
\end{equation}
where
\begin{equation}
  w_i\geq 0,\qquad \sum_{i=1}^n w_i=1,\qquad \eta>0.
\end{equation}
The actual privacy loss allocated to participant \(i\) is
\begin{equation}
  \epsilon_i^{\mathrm{actual}}=\eta w_i.
\end{equation}
It is useful to define
\begin{equation}
  y_i=\eta w_i,\qquad S=\sum_i y_i=\eta.
  \label{eq:yS}
\end{equation}
Then
\begin{equation}
  w_i=\frac{y_i}{S},\qquad \sum_i y_i=S,
\end{equation}
and the Laplace noise has scale \(1/S\) and variance \(2/S^2\).
Thus \(y_i\) is the actual privacy loss assigned to participant \(i\), while
\(S\) is the total privacy-loss parameter.

\subsection{Original statistical setting}

In the original common-population setting, observations are drawn from a
common distribution with mean \(\theta\) and variance \(\sigma^2\). For
independent observations, the variance of the weighted data component is
\(\sigma^2\sum_i w_i^2\). Since the estimator is unbiased in that setting,
its mean-squared error is
\begin{equation}
  \MSE
  =\sigma^2\sum_i w_i^2+\frac{2}{\eta^2}.
\end{equation}
In \(y,S\) variables this becomes
\begin{equation}
  \MSE
  =\frac{\sigma^2\sum_i y_i^2+2}{S^2}.
\end{equation}

\subsection{Virtual costs and payment identity}

Let participant \(i\) have private cost type \(c_i\), with distribution
\(F_i\) and density \(f_i\). Under the usual regularity assumptions, the
virtual cost is
\begin{equation}
  \psi_i(c_i)=c_i+\frac{F_i(c_i)}{f_i(c_i)}.
  \label{eq:virtual-cost}
\end{equation}
The source paper transforms the platform's expected payment objective into
an objective involving virtual costs.

For the participant utility and outside option used in the original model,
the payment identity has the form
\begin{equation}
  t_i(c_i)
  =h_i(c_i)-\operatorname{var}
   +c_i\epsilon_i(c_i)
   +\int_{c_i}^{\infty}\epsilon_i(z)\,dz,
  \label{eq:payment-identity}
\end{equation}
where \(h_i(c_i)\) is the relevant interim expected estimation loss and
\(\operatorname{var}\) denotes the participant's outside-option benchmark
in the source model. The precise interpretation of \(h_i\) follows the
source paper's utility specification.

The associated incentive-compatibility condition is that the interim
privacy allocation is weakly decreasing in the participant's own cost:
\begin{equation}
  c_i'\geq c_i
  \quad\Longrightarrow\quad
  \epsilon_i(c_i')\leq \epsilon_i(c_i).
\end{equation}
The payment identity itself is unchanged by the fairness extension below,
because fairness is introduced as a platform-level objective term and
participants' utilities are not changed.

\subsection{Original transformed optimization}

After substituting the payment identity and taking expectations, the
source paper obtains an objective of the form
\begin{equation}
  \E\left[
    (n+1)\MSE+\sum_i\psi_i(c_i)\epsilon_i(c)
  \right]-n\operatorname{var}.
\end{equation}
The additive constant \(-n\operatorname{var}\) does not affect the
allocation optimization. In the original common-distribution setting,
the pointwise optimization is
\begin{equation}
  \min_{y_i\geq0}
  \frac{n+1}{(\sum_i y_i)^2}
  \left(2+\sigma^2\sum_i y_i^2\right)
  +\sum_i\psi_i y_i.
  \label{eq:original-objective}
\end{equation}
This objective is not jointly convex in \(y\), because
\(S=\sum_i y_i\) is endogenous.

\subsection{How the paper obtains its threshold sweep}

The paper fixes \(S=\sum_i y_i\), solves the resulting convex inner
optimization, and uses the KKT conditions:
\begin{equation}
  y_i=
  \left(
    \frac{(\ell-\psi_i)S^2}
    {2(n+1)\sigma^2}
  \right)_+.
  \label{eq:paper-kkt}
\end{equation}
The normalization \(\sum_i y_i=S\) gives
\begin{equation}
  S(\ell)
  =
  \frac{2(n+1)\sigma^2}
       {\sum_i(\ell-\psi_i)_+}.
  \label{eq:paper-S}
\end{equation}
Thus the paper eliminates \(S\) in favour of one scalar threshold
\(\ell\). The active set changes only when \(\ell\) crosses a sorted virtual
cost. The paper's Algorithm 1 sweeps these intervals and obtains an
\(O(n\log n)\) algorithm (sorting plus a linear sweep).

The important structural feature is that every participant has the same
threshold \(\ell\), shifted only by their virtual cost. The fairness
extension creates group-specific thresholds, so the one-threshold
derivation does not transfer verbatim.

\section{Two-group statistical model}

\subsection{Group distributions and target}

Partition participants into two nonempty groups \(A\) and \(B\). Participant
\(i\) belongs to group \(g(i)\in\{A,B\}\). Assume
\begin{equation}
  \E[X_i]=\mu_{g(i)},\qquad
  \operatorname{Var}(X_i)=\sigma^2,
  \label{eq:group-moments}
\end{equation}
with a common within-group variance \(\sigma^2\). The common-variance
assumption is deliberate and is used throughout the current derivation.

Let \(p_A,p_B>0\), \(p_A+p_B=1\), and define the target population mean
\begin{equation}
  \theta=p_A\mu_A+p_B\mu_B.
  \label{eq:population-mean}
\end{equation}
For weights \(w\), define the estimator's mean
\begin{equation}
  m(w)=\sum_i w_i\mu_{g(i)}.
\end{equation}
The estimator's variance is
\begin{equation}
  V(w,S)=\sigma^2\sum_iw_i^2+\frac{2}{S^2}.
\end{equation}
Its group-conditional MSEs are
\begin{equation}
  \MSE_A=V(w,S)+(m(w)-\mu_A)^2,
  \qquad
  \MSE_B=V(w,S)+(m(w)-\mu_B)^2.
  \label{eq:group-MSE}
\end{equation}

\subsection{MSE disparity and fairness penalty}

Subtracting the group MSEs cancels the common variance term:
\begin{align}
  D(w)
  &:=\MSE_A-\MSE_B\\
  &=(m-\mu_A)^2-(m-\mu_B)^2\\
  &=(\mu_B-\mu_A)(2m-\mu_A-\mu_B).
  \label{eq:disparity}
\end{align}
Write
\begin{equation}
  \Delta=\mu_B-\mu_A.
\end{equation}
The platform-level fairness penalty is
\begin{equation}
  \mathcal F(w)=\lambda D(w)^2
  =\lambda\Delta^2(2m(w)-\mu_A-\mu_B)^2,
  \qquad \lambda\geq0.
  \label{eq:fairness}
\end{equation}
The population MSE is
\begin{equation}
  \MSE_{\mathrm{pop}}
  =p_A\MSE_A+p_B\MSE_B
  =V(w,S)+(m(w)-\theta)^2.
  \label{eq:pop-MSE}
\end{equation}

\paragraph{Interpretation.}
The fairness term penalizes the squared difference between the two
groups' conditional MSEs. It is not the same as penalizing the difference
between group means, nor is it the same as penalizing the variance of
individual losses. Here it is a deterministic function of the weighted
estimator mean \(m(w)\).

\subsection{Platform-level fairness only}

The fairness penalty is included in the platform's objective. It is not
included in participants' utilities. Therefore:
\begin{itemize}
  \item the participants' utility model and outside option remain those of
        the original paper;
  \item the payment identity remains unchanged;
  \item the original monotonicity condition on interim privacy allocations
        remains the relevant IC condition;
  \item fairness changes the platform's allocation objective, not the
        participant's direct incentive.
\end{itemize}

In particular, the transformed objective contains
\(\lambda\mathcal F\), not \((n+1)\lambda\mathcal F\). The factor \(n+1\)
multiplies the population MSE because of the original payment-identity
transformation; it does not multiply a new platform-only fairness term.

\section{Fairness-augmented mechanism-design objective}

Substituting the two-group population MSE and fairness penalty into the
platform's transformed objective gives
\begin{equation}
  \begin{aligned}
  \min_{w,S}\quad J(w,S)
  ={}&(n+1)\left[
       \sigma^2\sum_iw_i^2+\frac{2}{S^2}
       +(m(w)-\theta)^2
       \right]\\
     &+\lambda\Delta^2(2m(w)-\mu_A-\mu_B)^2
       +S\sum_i\psi_iw_i,
  \end{aligned}
  \label{eq:full-w-objective}
\end{equation}
subject to
\begin{equation}
  w_i\geq0,\qquad \sum_iw_i=1,\qquad S>0.
\end{equation}
Here \(\psi_i=\psi_i(c_i)\) is the virtual cost at the reported type
profile. The pointwise optimization is performed for each profile of
virtual costs.

Using \(y_i=Sw_i\), \(S=\sum_i y_i\), the same objective is
\begin{equation}
  \begin{aligned}
  J(y,S)
  ={}&(n+1)\left[
      \frac{\sigma^2\sum_i y_i^2}{S^2}
      +\frac{2}{S^2}
      +\left(\frac{\sum_i\mu_{g(i)}y_i}{S}-\theta\right)^2
      \right]\\
    &+\lambda\Delta^2
      \left(
      2\frac{\sum_i\mu_{g(i)}y_i}{S}
      -\mu_A-\mu_B
      \right)^2
      +\sum_i\psi_i y_i,
  \end{aligned}
  \label{eq:full-y-objective}
\end{equation}
with \(y_i\geq0\) and \(S=\sum_i y_i>0\).

\section{Fixed-\texorpdfstring{$S$}{S} optimization and KKT conditions}

\subsection{Inner problem}

Fix \(S>0\). Define
\begin{equation}
  a_i=\frac{\mu_{g(i)}}{S},
  \qquad
  m=a^\top y,
  \qquad
  \kappa=\frac{2(n+1)\sigma^2}{S^2}.
\end{equation}
The fixed-\(S\) problem is
\begin{equation}
  \begin{aligned}
  \min_{y_i\geq0}\quad&
  \frac{2(n+1)}{S^2}
  +\frac{(n+1)\sigma^2}{S^2}\sum_i y_i^2
  +(n+1)(a^\top y-\theta)^2\\
  &+\sum_i\psi_i y_i
  +\lambda\Delta^2
      (2a^\top y-\mu_A-\mu_B)^2\\
  \text{subject to}\quad&\sum_i y_i=S.
  \end{aligned}
  \label{eq:fixed-S}
\end{equation}

Define
\begin{equation}
  L=2(n+1)+8\lambda\Delta^2,
  \qquad
  q_0=-2(n+1)\theta-4\lambda\Delta^2(\mu_A+\mu_B),
  \label{eq:Lq0}
\end{equation}
and
\begin{equation}
  q(m)=Lm+q_0.
  \label{eq:q}
\end{equation}

The Hessian of the objective with respect to \(y\) is
\begin{equation}
  H_y
  =
  \frac{2(n+1)\sigma^2}{S^2}I
  +\left(2(n+1)+8\lambda\Delta^2\right)aa^\top.
  \label{eq:hessian}
\end{equation}
For every nonzero \(v\),
\begin{equation}
  v^\top H_yv
  =
  \frac{2(n+1)\sigma^2}{S^2}\|v\|^2
  +L(a^\top v)^2>0.
\end{equation}
Thus the fixed-\(S\) problem is strictly convex and has a unique global
minimizer.

\subsection{KKT allocation rule}

Let \(\gamma\) be the multiplier for \(\sum_i y_i=S\), and put
\(z=-\gamma\). The KKT conditions give
\begin{equation}
  \boxed{
  y_i
  =
  \frac{(z-\psi_i-q(m)a_i)_+}{\kappa}.
  }
  \label{eq:kkt-y}
\end{equation}
The scalar \(z\) is chosen to satisfy the normalization
\(\sum_i y_i=S\), while \(q\) must be consistent with \(m\).

The fairness term makes the allocation depend on both the participant's
virtual cost and their group mean. This is why one common threshold no
longer describes all participants.

\subsection{Monotonicity in own virtual cost and IC}

For fixed \(S\), the objective can be written as
\[
G(y)+\sum_i\psi_i y_i,
\qquad
\sum_i y_i=S,\quad y_i\geq0,
\]
where \(G\) is independent of the reported costs. Comparing global
minimizers at two values \(\psi_i'<\!\!{}^{\,c}\) is most cleanly done
using the standard two-optimum argument. If only participant \(i\)'s
virtual cost rises from \(\psi_i\) to \(\psi_i'>\psi_i\), optimality at
the two profiles implies
\begin{equation}
  (\psi_i'-\psi_i)(y_i'-y_i)\leq0,
\end{equation}
so
\begin{equation}
  y_i'\leq y_i.
\end{equation}
Because \(S\) is fixed, \(\epsilon_i=y_i\), and privacy allocation is
weakly decreasing in own virtual cost. If \(\psi_i(c_i)\) is
nondecreasing in \(c_i\), this gives the required monotonicity in own
cost.

The argument uses global optimality and a report-independent feasible
set; it does not require convexity. For the full problem where \(S\) is
also optimized, the same comparison applies if the feasible set and all
other terms remain independent of the participant's report and the
global optimum is selected. The fixed-\(S\) strict-convexity result alone
does not by itself prove every Bayesian IC detail; the interim allocation
and tie-breaking assumptions should be stated in a final mechanism proof.

\section{Eliminating the total privacy parameter}

\subsection{Why the fixed-\texorpdfstring{$S$}{S} sweep is not yet enough}

For a fixed \(S\), the inner problem is convex and can be solved by KKT.
But the platform must also choose \(S\). Solving the inner problem for
one \(S\) does not solve the original joint optimization.

The objective in \((w,S)\) has the form
\begin{equation}
  J(w,S)=Q(w)+\frac{2(n+1)}{S^2}
  +S\sum_i\psi_iw_i,
  \label{eq:Q-form}
\end{equation}
where
\begin{equation}
  Q(w)=(n+1)\left[
    \sigma^2\sum_iw_i^2+(m(w)-\theta)^2
  \right]
  +\lambda\Delta^2(2m(w)-\mu_A-\mu_B)^2.
\end{equation}

For fixed \(w\), if \(\bar\psi(w)=\sum_i\psi_iw_i>0\), minimizing over
\(S\) gives
\begin{equation}
  S^*(w)=
  \left(\frac{4(n+1)}{\bar\psi(w)}\right)^{1/3}.
  \label{eq:S-fixed-w}
\end{equation}
Substitution gives the reduced problem
\begin{equation}
  \min_{w\in\Delta_n}
  Q(w)+\frac32[4(n+1)]^{1/3}
  \left(\sum_i\psi_iw_i\right)^{2/3}.
  \label{eq:reduced-w}
\end{equation}
The power \(2/3\) is concave, so this reduced objective is not
automatically convex. This reparameterization is useful for understanding
the outer problem, but it does not directly yield a convex optimization
algorithm.

\subsection{Envelope condition}

Let
\[
\phi(S)=\min_{\substack{y_i\geq0\\\sum_i y_i=S}}J(y,S).
\]
The envelope theorem gives, at differentiability points,
\begin{equation}
  \phi'(S)
  =
  \sum_i\psi_iw_i^*(S)-\frac{4(n+1)}{S^3}.
  \label{eq:envelope}
\end{equation}
An interior optimum therefore satisfies
\begin{equation}
  \boxed{
  \sum_i\psi_iw_i^*(S)
  =\frac{4(n+1)}{S^3}.
  }
  \label{eq:outer-stationarity}
\end{equation}
No uniqueness of this outer root is assumed. In particular, bisection
would require a separate proof of a suitable monotonicity or sign-change
property.

\subsection{Two-variable reparameterization}

The fixed-\(S\) KKT rule is
\[
y_i=\frac{(z-\psi_i-q a_i)_+}{\kappa},
\quad
\kappa=\frac{K}{S^2},
\quad
a_i=\frac{\mu_{g(i)}}{S},
\quad
K=2(n+1)\sigma^2.
\]
Set
\[
r=\frac qS,
\qquad
u_i(z,r)=(z-\psi_i-r\mu_{g(i)})_+,
\qquad
H(z,r)=\sum_i u_i(z,r).
\]
Then
\begin{equation}
  y_i=\frac{S^2}{K}u_i,
  \qquad
  S=\frac K H.
  \label{eq:S-H}
\end{equation}
Define
\begin{equation}
  m(z,r)=
  \frac{\sum_i\mu_{g(i)}u_i(z,r)}{H(z,r)}.
\end{equation}
The consistency relation \(q=Lm+q_0\), together with \(q=rS\),
becomes
\begin{equation}
  \boxed{rK=(Lm(z,r)+q_0)H(z,r).}
  \label{eq:consistency-zr}
\end{equation}
The outer stationarity condition becomes
\begin{equation}
  \boxed{
  \frac{\sum_i\psi_i u_i(z,r)}{H(z,r)}
  =\frac{4(n+1)H(z,r)^3}{K^3}.
  }
  \label{eq:outer-zr}
\end{equation}
This is a system of two equations in \(z,r\). Within a fixed active set,
the positive-part expressions become affine, so the first equation is
linear and the second is quartic after eliminating one variable.

\section{A monotone two-threshold path}

\subsection{Threshold coordinates}

Define group-specific thresholds
\begin{equation}
  \tau_A=z-r\mu_A,
  \qquad
  \tau_B=z-r\mu_B.
  \label{eq:thresholds}
\end{equation}
Then
\begin{equation}
  r=\frac{\tau_A-\tau_B}{\Delta},
\end{equation}
and
\begin{equation}
  u_i=(\tau_{g(i)}-\psi_i)_+.
\end{equation}
Define
\begin{equation}
  H_A(\tau_A)=\sum_{i\in A}(\tau_A-\psi_i)_+,
  \qquad
  H_B(\tau_B)=\sum_{i\in B}(\tau_B-\psi_i)_+.
\end{equation}
Thus \(H=H_A+H_B\).

The consistency equation \eqref{eq:consistency-zr} is equivalent to
\begin{equation}
  \boxed{
  F(\tau_A,\tau_B)
  =
  \frac K\Delta(\tau_A-\tau_B)
  -c_AH_A(\tau_A)-c_BH_B(\tau_B)=0,
  }
  \label{eq:F}
\end{equation}
where
\begin{equation}
  c_A=L\mu_A+q_0,
  \qquad
  c_B=L\mu_B+q_0.
\end{equation}

\subsection{Sign structure}

Using \(\theta=p_A\mu_A+p_B\mu_B\), with both proportions strictly
positive, we have
\begin{align}
  c_A
  &=2(n+1)(\mu_A-\theta)-4\lambda\Delta^3<0,
  \label{eq:cA-sign}\\
  c_B
  &=2(n+1)(\mu_B-\theta)+4\lambda\Delta^3>0.
  \label{eq:cB-sign}
\end{align}
These signs are central to the sweep result.

\begin{proposition}[Unique monotone consistency path]
Suppose \(\mu_A<\mu_B\), \(\theta\) is strictly between the group means,
\(\sigma^2>0\), and \(\lambda\geq0\). Then, for every
\(\tau_A\in\R\), there is exactly one \(\tau_B\in\R\) satisfying
\eqref{eq:F}. The resulting function
\(\tau_B=g(\tau_A)\) is continuous, strictly increasing, and
piecewise linear.
\end{proposition}

\begin{proof}
For fixed \(\tau_A\), the function \(F\) is continuous in \(\tau_B\).
Its derivative wherever it exists is
\[
  \frac{\partial F}{\partial\tau_B}
  =-\frac K\Delta-c_B k_B<0,
\]
where \(k_B\geq0\) is the number of active B participants. Thus it is
strictly decreasing.

As \(\tau_B\to-\infty\), \(H_B=0\) and
\(F\to+\infty\). As \(\tau_B\to+\infty\), \(H_B\to+\infty\) and,
because \(c_B>0\), \(F\to-\infty\). Hence there is exactly one root.

Similarly, wherever derivatives exist,
\[
  \frac{\partial F}{\partial\tau_A}
  =\frac K\Delta-c_Ak_A>0.
\]
Implicit differentiation therefore gives
\[
  \frac{d\tau_B}{d\tau_A}
  =-\frac{F_{\tau_A}}{F_{\tau_B}}>0.
\]
Continuity follows from uniqueness and the continuous, strictly
monotone dependence of \(F\) on \(\tau_B\). On any region where the
active sets are fixed, \(H_A,H_B\) are affine, so \(F=0\) is an affine
relation between the thresholds. Thus the path is piecewise linear.
\end{proof}

\begin{remark}[No participant leaves]
Both thresholds increase along the path. Since a participant is active
exactly when \(\tau_g>\psi_i\), participants enter in increasing
virtual-cost order and never leave. Starting below all virtual costs,
the path eventually activates all participants. There are at most \(n\)
entry events and at most \(n+1\) segments.
\end{remark}

\subsection{Explicit segment equation}

Within a segment, let \(k_A,k_B\) be the active counts and let
\[
P_A=\sum_{i\in\mathcal A_A}\psi_i,\qquad
P_B=\sum_{i\in\mathcal A_B}\psi_i.
\]
Then
\[
H_A=k_A\tau_A-P_A,\qquad
H_B=k_B\tau_B-P_B.
\]
Equation \eqref{eq:F} becomes
\begin{equation}
  A_0\tau_A+B_0\tau_B+C_0=0,
  \label{eq:segment-line}
\end{equation}
where
\begin{align}
  A_0&=\frac K\Delta-c_Ak_A>0,\\
  B_0&=-\frac K\Delta-c_Bk_B<0,\\
  C_0&=c_AP_A+c_BP_B.
\end{align}
Consequently,
\begin{equation}
  \boxed{
  \tau_B=s\tau_A+d,\qquad
  s=-\frac{A_0}{B_0}>0,\qquad
  d=-\frac{C_0}{B_0}.
  }
  \label{eq:segment-path}
\end{equation}

\subsection{Event-driven enumeration}

Sort virtual costs separately:
\[
\psi_{A,(1)}\leq\cdots\leq\psi_{A,(n_A)},
\qquad
\psi_{B,(1)}\leq\cdots\leq\psi_{B,(n_B)}.
\]
Because both thresholds increase, the active participants in each
group are always a prefix of its sorted list.

For a current segment, the next events are:
\begin{itemize}
  \item A-entry:
  \[
    \tau_A=\psi_{A,(k_A+1)}.
  \]
  \item B-entry, using \(\tau_B=s\tau_A+d\):
  \[
    \tau_A=
    \frac{\psi_{B,(k_B+1)}-d}{s}.
  \]
\end{itemize}
Choose the smallest event parameter strictly ahead of the current
position. If the two events coincide, update both groups simultaneously.
Recompute the segment coefficients and continue.

The path begins with both groups inactive and
\(\tau_A=\tau_B\). Its first event is the smallest virtual cost across
both groups. The path eventually reaches the region in which all
participants are active.

\section{Quartic equation on each segment}

\subsection{Prefix-sum expressions}

For a fixed active set, define
\begin{equation}
  R_A=\sum_{i\in\mathcal A_A}\psi_i^2,
  \qquad
  R_B=\sum_{i\in\mathcal A_B}\psi_i^2.
\end{equation}
The two quantities in the outer stationarity equation are
\begin{align}
  H
  &=k_A\tau_A+k_B\tau_B-P_A-P_B,\\
  U
  &=P_A\tau_A+P_B\tau_B-R_A-R_B,
  \label{eq:HU}
\end{align}
where
\[
  U=\sum_i\psi_i u_i.
\]
Using \(\tau_B=s\tau_A+d\), both are affine in \(\tau_A\):
\begin{equation}
  H=h_0+h_1\tau_A,\qquad
  U=u_0+u_1\tau_A,
\end{equation}
with
\begin{align}
  h_0&=k_Bd-P_A-P_B,
  &h_1&=k_A+k_Bs,\\
  u_0&=P_Bd-R_A-R_B,
  &u_1&=P_A+P_Bs.
\end{align}

\subsection{Quartic stationarity equation}

The outer stationarity equation \eqref{eq:outer-zr} is
\[
  K^3U=4(n+1)H^4.
\]
Substituting the affine expressions gives
\begin{equation}
  \boxed{
  4(n+1)(h_0+h_1\tau_A)^4
  -K^3(u_0+u_1\tau_A)=0.
  }
  \label{eq:quartic}
\end{equation}
This polynomial has degree at most four.

For each segment:
\begin{enumerate}
  \item Solve the polynomial for all real roots.
  \item Retain only roots in the segment's threshold interval.
  \item Require \(H>0\), so that \(S=K/H>0\).
  \item Check the active and inactive inequalities.
  \item Reconstruct
  \[
    S=\frac K H,\qquad
    w_i=\frac{(\tau_{g(i)}-\psi_i)_+}{H},
    \qquad
    y_i=Sw_i.
  \]
\end{enumerate}
All valid roots must be retained; there is no established result that
there is at most one root per segment.

\section{Global-optimality argument}

\subsection{Existence of a finite minimizer}

For the finite-minimizer argument, assume
\begin{equation}
  \psi_i\geq\psi_{\min}>0
  \qquad\text{for every }i.
  \label{eq:positive-psi}
\end{equation}
The objective can be written
\[
  J(w,S)=Q(w)+\frac{2(n+1)}{S^2}
       +S\sum_i\psi_iw_i,
\]
where \(Q(w)\geq0\). Since \(w\) is in the probability simplex,
\[
  J(w,S)\geq
  \frac{2(n+1)}{S^2}+S\psi_{\min}.
\]
The right-hand side tends to \(+\infty\) both as \(S\to0^+\) and
as \(S\to+\infty\), uniformly over \(w\). The simplex is compact
and the objective is continuous, so a global minimizer exists at a
finite \(S^*>0\).

\subsection{Why the fixed-\texorpdfstring{$S$}{S} KKT solution is globally optimal}

For each fixed \(S\), the objective is strictly convex in \(w\)
(equivalently in \(y\) subject to \(\sum_i y_i=S\)). Therefore its KKT
conditions are necessary and sufficient, and its minimizer is unique.

Every point on the consistency path with \(H>0\) defines
\[
  S=K/H,\qquad
  w_i=u_i/H.
\]
The consistency equation ensures that the \(q\) used in the KKT rule
equals \(Lm+q_0\). The normalization is satisfied because
\[
  \sum_i y_i
  =\frac{S^2}{K}H
  =S.
\]
Thus each such point gives the unique globally optimal allocation for
its corresponding fixed \(S\).

\subsection{Why the global optimum is among the quartic candidates}

Let \((w^*,S^*)\) be a global minimizer. Under
\eqref{eq:positive-psi}, it exists at a finite, positive \(S^*\).
The unique fixed-\(S^*\) minimizer must satisfy the KKT conditions.
Moreover, since \(S^*\) is an interior optimum of the outer problem,
the envelope derivative vanishes:
\[
  \sum_i\psi_iw_i^*-\frac{4(n+1)}{(S^*)^3}=0.
\]
Hence its thresholds lie on the consistency path and satisfy the
quartic stationarity equation on their active-set segment (or at a
segment endpoint).

Conversely, every valid quartic root on the path gives a feasible
allocation satisfying the fixed-\(S\) KKT conditions and outer
stationarity. Evaluating the original objective at all such candidates
and selecting the smallest value therefore recovers the global
minimizer.

\begin{proposition}[Global candidate characterization]
Under the assumptions of the monotone-path proposition and
\(\psi_{\min}>0\), every global minimizer of
\eqref{eq:full-w-objective} is represented by a valid real root of
\eqref{eq:quartic} on one of the at most \(n+1\) active-set segments,
including shared segment endpoints. Consequently, evaluating all valid
candidates and selecting the one with the smallest original objective
value gives the global optimum.
\end{proposition}

\begin{proof}
Existence follows from coercivity. Fixed-\(S\) strict convexity makes
the KKT allocation unique and globally optimal. The outer first-order
condition is necessary at the finite interior global minimizer.
The threshold-path proposition enumerates every possible active set
along the unique consistency path, and the quartic equation is
equivalent to the outer first-order condition on each segment.
Therefore the global minimizer is included in the candidate set.
Taking the smallest objective value among feasible candidates cannot
produce a value below the global minimum, and the candidate set contains
a global minimizer. Hence the smallest candidate value equals the
global minimum.
\end{proof}

\subsection{Zero virtual costs and the limiting boundary}

If all virtual costs are nonnegative but some are zero, the finite-\(S\)
coercivity argument may fail. As \(S\to\infty\), bounded objective
value forces all limiting weight to be placed on zero-virtual-cost
participants. Define
\[
  Z_0=\{i:\psi_i=0\}.
\]
If \(Z_0\neq\varnothing\), the limiting boundary value is
\begin{equation}
  \boxed{
  J_\infty
  =
  \min_{\substack{w_i\geq0,\ \sum_iw_i=1\\
                  w_i=0\text{ for }i\notin Z_0}}
  Q(w).
  }
  \label{eq:infinite-boundary}
\end{equation}
This value must be compared with all finite stationary candidates.
The infimum may occur only in the limit \(S\to\infty\), rather than at
a finite mechanism. If \(Z_0\) is empty, nonnegative virtual costs
imply \(\psi_{\min}>0\), and the finite-minimizer argument applies.

If negative virtual costs are possible, the analysis changes: the
objective can become unbounded below as \(S\to\infty\) along a
negative-cost participant. The current theorem therefore assumes
nonnegative virtual costs.

\section{Algorithm and complexity}

\subsection{Event-driven algorithm}

\begin{enumerate}
  \item Sort virtual costs separately within groups A and B.
  \item Initialize \(\tau_A=\tau_B\) below all virtual costs, with both
        groups inactive.
  \item Determine the first participant-entry event.
  \item For each active-set segment:
  \begin{enumerate}
    \item Compute \(k_A,k_B,P_A,P_B,R_A,R_B\).
    \item Compute the affine threshold relation
          \(\tau_B=s\tau_A+d\).
    \item Compute the next A-entry and B-entry events.
    \item Construct the quartic \eqref{eq:quartic}.
    \item Solve it and retain all valid roots in the current segment.
    \item Advance to the next event; process simultaneous events together.
  \end{enumerate}
  \item Include segment endpoints and deduplicate candidates.
  \item If zero virtual costs exist, compute the limiting value
        \eqref{eq:infinite-boundary}.
  \item Evaluate the original objective at every candidate and boundary
        value; return the smallest.
\end{enumerate}

\subsection{Complexity}

Sorting the two groups costs
\[
  O(n_A\log n_A+n_B\log n_B)=O(n\log n).
\]
Prefix sums of virtual costs and squared virtual costs can be built in
\(O(n)\). Each participant enters once, so there are at most \(n\)
events and \(n+1\) segments. A constant-degree quartic is solved a
constant number of times per segment in the arithmetic-operation model.
Candidate reconstruction and objective evaluation are linear in the
number of participants if performed naively, giving a total of
\[
  \boxed{O(n\log n)}
\]
arithmetic operations, with \(O(n)\) memory.

This is an algebraic complexity claim. A numerical implementation still
needs a robust quartic root-isolation routine and careful treatment of
near-coincident events, repeated roots, and degenerate polynomials.

\section{Earlier fixed-\texorpdfstring{$S$}{S} active-set sweep}

Before deriving the joint \((w,S)\) sweep, a separate active-set
algorithm was developed for the fixed-\(S\) inner problem. It is useful
both as an independent implementation route and as a way to validate
the joint algorithm.

For fixed \(S\), write
\[
  y_i(q)=\frac{(z(q)-qa_i-\psi_i)_+}{\kappa},
  \qquad
  \sum_i(z-qa_i-\psi_i)_+=\kappa S.
\]
The consistency condition is
\[
  q=Lm(q)+q_0,
  \qquad
  m(q)=\sum_i a_i y_i(q).
\]
Define
\[
  h(q)=q-Lm(q)-q_0.
\]
Within an active-set interval, \(m(q)\) is affine. In the two-group
case, let the group coefficients be ordered \(a_1<a_2\). If the active
participants are prefixes of sizes \(k_1,k_2\), define
\[
  P_g=\sum_{j=1}^{k_g}\psi_{g,(j)},\quad
  k=k_1+k_2,\quad
  A_1=k_1a_1+k_2a_2.
\]
Normalization gives
\[
  z(q)=\frac{\kappa S+P_1+P_2+qA_1}{k}.
\]
The group thresholds are
\[
  t_g(q)=z(q)-qa_g,
\]
with slopes
\[
  t_g'(q)=A_1/k-a_g.
\]
For \(a_1<a_2\), the group-1 threshold is nondecreasing and the
group-2 threshold is nonincreasing. Hence group 1 participants enter
in ascending virtual-cost order, while group 2 participants leave from
the highest virtual-cost end. This produces a different event structure
from the joint-\(S\) sweep, where both groups' thresholds increase.

For a fixed active set, define
\[
  A_2=k_1a_1^2+k_2a_2^2,
\]
\[
  B=\frac{A_1(\kappa S+P_1+P_2)}{k}
    -(a_1P_1+a_2P_2),
\qquad
  C=\frac{A_1^2}{k}-A_2\leq0.
\]
Then
\[
  m(q)=\frac{B+Cq}{\kappa},
\]
and the consistency root is
\[
  q^*=\frac{\kappa q_0+LB}{\kappa-LC}.
\]
The event sweep can locate the interval containing this root, after
which \(z\) and \(y_i\) are recovered.

This fixed-\(S\) method is a candidate \(O(n\log n)\) inner solver.
Its event logic, especially initialization, simultaneous events, and
the cases in which one group is empty, should be independently tested.
It does not optimize over \(S\); the joint two-threshold method above
does.

\section{Mechanism-design implications and limitations}

\subsection{IC implications}

The fairness term is platform-only, so the participant utility and
payment identity are unchanged. Under independent types, regular
(non-decreasing) virtual costs, the original utility specification,
and a pointwise globally optimal allocation rule, monotonicity in own
virtual cost is supported by the two-optimum comparison. This is the
key reason the pointwise optimization can potentially be used in the
Bayesian mechanism.

A final mechanism theorem should still state carefully:
\begin{itemize}
  \item the type distributions and independence assumptions;
  \item regularity of virtual costs;
  \item whether the pointwise optimizer is unique (it is unique for fixed
        \(S\), while the full joint problem may have multiple global
        minimizers);
  \item how ties between global minimizers are resolved;
  \item the exact interim allocation and payment definitions.
\end{itemize}
The algorithmic optimization result alone is not a complete Bayesian
IC/IR proof.

\subsection{Estimator optimality versus optimization optimality}

There are two distinct claims.

\paragraph{Claim 1: global optimization within the chosen estimator class.}
The algorithm minimizes the fairness-augmented platform objective over
all linear weighted estimators with Laplace noise of the form
\eqref{eq:estimator}, subject to the stated allocation structure.
This is an optimization statement about a restricted mechanism class.

\paragraph{Claim 2: near-optimality among all DP estimators.}
This would compare the chosen estimator with every admissible central-DP
mechanism, including nonlinear estimators and mechanisms using other
noise distributions or data-dependent procedures. The optimization
algorithm does not establish this claim.

A possible formal target is a heterogeneous minimax risk. Let
\[
  \mathcal P_{\mathrm{het}}
  =\mathcal P_A^{n_A}\times\mathcal P_B^{n_B}
\]
be a specified class of group distributions, and let
\(\mathcal Q_{\mathrm{DP}}(\epsilon)\) be the class of all central-DP
estimators satisfying the relevant privacy constraints. Define
\begin{equation}
  R_{\mathrm{all}}^*(\epsilon)
  =
  \inf_{\mathcal M\in\mathcal Q_{\mathrm{DP}}(\epsilon)}
  \sup_{(P_A,P_B)\in\mathcal P_{\mathrm{het}}}
  \E\left[
    \bigl(\mathcal M(X)-\theta(P_A,P_B)\bigr)^2
  \right].
\end{equation}
One would then seek an upper bound for the proposed estimator and a
matching or comparable lower bound for \(R_{\mathrm{all}}^*\).

The original paper's near-optimality result is established for its
original statistical setting (iid observations from a common
distribution and its specified minimax class). It does not automatically
transfer to two heterogeneous group distributions, a population-weighted
target, and a fairness-augmented objective. A new lower-bound argument
and a precise definition of the fairness-aware risk criterion would be
needed.

\subsection{What the current result does and does not establish}

\begin{center}
\begin{tabular}{p{0.43\textwidth}p{0.48\textwidth}}
\toprule
\textbf{Established by the current derivation} &
\textbf{Not established yet}\\
\midrule
Fixed-\(S\) objective is strictly convex and has a unique allocation. &
Near-optimality among all central-DP estimators.\\
The two-threshold consistency path is unique, increasing, and
piecewise linear under the stated assumptions. &
A statistical minimax lower bound for the heterogeneous two-group model.\\
There are at most \(n\) participant-entry events. &
A complete end-to-end numerical stability guarantee.\\
The outer stationarity equation is quartic on each segment. &
A complete mechanism theorem without specifying tie-breaking and all
Bayesian assumptions.\\
With positive virtual costs, the valid stationary candidates contain
the global optimum. &
Extension to unequal within-group variances or more than two groups.\\
\bottomrule
\end{tabular}
\end{center}

\section{Compact conceptual summary}

The research extension can be understood as follows.

\begin{enumerate}
  \item The original paper chooses privacy allocations to balance
        estimation error and participant virtual costs.
  \item We introduce two groups with different means and penalize the
        disparity between their conditional MSEs.
  \item This fairness term changes the KKT allocation rule: participants
        now face group-dependent thresholds.
  \item At fixed total privacy loss \(S\), the allocation problem remains
        strictly convex.
  \item To optimize \(S\) as well, we combine KKT consistency with the
        outer envelope condition.
  \item The consistency equation couples the two thresholds. The
        population-mean assumption gives opposite signs for the two
        group coefficients, making the threshold path unique and
        monotone.
  \item Each participant enters once, so there are only linearly many
        active-set segments.
  \item On each segment, the outer stationarity condition is quartic.
        Solving these quartics and comparing valid candidates gives a
        global optimizer under the positive-virtual-cost assumption.
  \item This establishes efficient optimization within the
        linear-Laplace class. It does not establish minimax
        near-optimality among all DP mechanisms.
\end{enumerate}

\end{document}
